% Figure: combined-cycle temperature-entropy schematic. The Brayton topping
% cycle sits at high temperature; the Rankine bottoming cycle sits below it,
% raised by the heat the gas-turbine exhaust hands to the heat-recovery steam
% generator. The two T-s loops are drawn schematically (not to scale) to show
% how the bottom cycle reuses the heat the top cycle would otherwise reject.
\begin{figure}[htbp]
\centering
\begin{tikzpicture}
\begin{axis}[
  width=12cm, height=8.5cm,
  xlabel={specific entropy $s$ (arbitrary units)},
  ylabel={temperature $T$ (K)},
  xmin=0, xmax=10, ymin=250, ymax=1650,
  axis lines=left,
  tick label style={font=\scriptsize}, label style={font=\small},
  clip=false,
  ytick={300,600,900,1200,1500},
]
% === Brayton topping cycle (schematic closed loop in the high-T band) ===
% 1 compressor inlet (low s, ~300), 2 compressor exit (~640), 3 turbine inlet
% (~1500, higher s), 4 turbine exit (~850, highest s), back to 1 along the
% rejection line. Drawn as a rounded quadrilateral.
\addplot[engred, very thick] coordinates {
  (1.0,300) (1.3,640) (4.6,1500) (6.5,850) (1.0,300)
};
\node[font=\scriptsize, engred, anchor=east] at (axis cs:1.0,300) {1};
\node[font=\scriptsize, engred, anchor=east] at (axis cs:1.3,640) {2};
\node[font=\scriptsize, engred, anchor=south] at (axis cs:4.6,1500) {3};
\node[font=\scriptsize, engred, anchor=west] at (axis cs:6.5,850) {4};
\node[font=\small, engred, anchor=west] at (axis cs:4.8,1250) {Brayton (gas turbine)};
% === Rankine bottoming cycle (schematic loop in the low-T band) ===
% a pump (~310), b boiler/HRSG heating up to ~800 at constant pressure,
% c superheated turbine inlet, d turbine exit two-phase near 310, back to a.
\addplot[engblue, very thick] coordinates {
  (3.2,310) (3.3,500) (5.0,800) (8.5,800) (8.7,310) (3.2,310)
};
\node[font=\scriptsize, engblue, anchor=east] at (axis cs:3.2,310) {a};
\node[font=\scriptsize, engblue, anchor=south] at (axis cs:5.0,800) {b};
\node[font=\scriptsize, engblue, anchor=south] at (axis cs:8.5,800) {c};
\node[font=\scriptsize, engblue, anchor=west] at (axis cs:8.7,310) {d};
\node[font=\small, engblue, anchor=west] at (axis cs:5.2,620) {Rankine (steam)};
% === The heat-transfer coupling: the Brayton rejection band overlaps the
% Rankine heat-addition band. Mark the shared temperature window. ===
\draw[enggray, dashed, thick] (axis cs:0,850) -- (axis cs:10,850);
\draw[enggray, dashed, thick] (axis cs:0,500) -- (axis cs:10,500);
\node[font=\scriptsize, enggray, anchor=west] at (axis cs:0.1,675)
  {HRSG transfer window};
\draw[-{Latex[length=2.2mm]}, enggray, very thick]
  (axis cs:7.0,840) -- (axis cs:7.0,510);
\end{axis}
\end{tikzpicture}
\caption[Combined-cycle temperature-entropy schematic]{Schematic
temperature-entropy view of a combined cycle, not drawn to scale. The Brayton
topping cycle (red) operates in the high-temperature band, with turbine inlet
near $1500\,$K; it rejects its exhaust heat over the temperature window between
the dashed lines. The heat-recovery steam generator hands that heat to the
Rankine bottoming cycle (blue), which raises and superheats steam and rejects
its own heat to the condenser near $310\,$K. The grey arrow marks the transfer
the HRSG performs: the heat the gas turbine would otherwise throw to the stack
becomes the heat input of the steam cycle. Because the two cycles occupy
different temperature bands, their efficiencies add far better than either
cycle could reach alone, which is the structural reason the combined cycle
clears $60\,\%$.}
\label{fig:vol04ch03:combined-cycle-ts}
\end{figure}
